\chapter{Methodology and Implementation}\label{ch:method}

\section{Research Design}\label{sec:design}

The study adopts a comparative empirical design in which four model pipelines are
evaluated across four publicly available credit-risk datasets under a common protocol.
Each dataset poses a binary default-classification task, which allows the models to be
compared while the credit product, the exact outcome definition, sample size, class
distribution, and temporal information vary.

Predictive and profit-based performance constitute the primary evaluation dimensions.
Attribution-ranking stability and cross-model agreement provide supplementary evidence
about the resulting explanation outputs. Cross-dataset comparisons, learning curves and
an out-of-time experiment examine how the observed findings vary across settings.

\section{Datasets and Prediction Tasks}\label{sec:datasets}

The four datasets differ substantially in sample size, class distribution and credit
product (Table~\ref{tab:datasets}). The first three are cross-sectional benchmarks of
increasing scale; Lending Club additionally enables an out-of-time evaluation.

\begin{table}[htbp]
\centering
\small
\begin{tabular}{lrccl}
\toprule
\textbf{Dataset} & \textbf{Samples} & \textbf{Predictors}
& \textbf{Default rate} & \textbf{Credit product} \\
\midrule
German Credit (Statlog) & 1{,}000       & 20 & 0.300 & Consumer instalment \\
Taiwan Default          & 30{,}000      & 23 & 0.221 & Revolving credit card \\
Home Credit             & 307{,}511     & 72 & 0.081 & General-purpose consumer \\
Lending Club            & 1{,}345{,}350 & 62 & 0.200 & Peer-to-peer lending \\
\bottomrule
\end{tabular}
\caption[The four credit datasets]{The four credit datasets, ordered by size.
Predictors denote retained source predictors. Default rate is the share of the positive
class.}
\label{tab:datasets}
\end{table}


Each dataset supplies its own published label rather than a horizon constructed here:
the good/bad credit classification of German Credit, delinquency in the following
billing cycle for the revolving card accounts of Taiwan, the documented
payment-difficulty target of Home Credit, and resolved loan status, fully paid against
charged off or defaulted, for Lending Club. Labels, horizons and populations therefore
differ, so the four are empirical contexts rather than interchangeable replications.

This thesis uses the original Statlog file \emph{German Credit} with its original attribute coding \citep{hofmannStatlogGermanCredit1994}. \emph{Taiwan Default} \citep{i-chengyehDefaultCreditCard2009} provides a
medium-sized benchmark for revolving credit card debt. It differs from the other datasets
because it represents behavioural scoring of existing card accounts rather than
application scoring: the predictors are six months of billing and repayment history, and
the target is delinquency in the month that follows. It is retained as an additional
credit context.
\emph{Home Credit Default Risk}\footnote{Kaggle competition dataset
\emph{Home Credit Default Risk} (2018),
\url{https://www.kaggle.com/c/home-credit-default-risk}, accessed 2~June~2026.}
represents general-purpose
consumer lending and exhibits the strongest class imbalance among the four
datasets.
\emph{Lending Club}\footnote{Kaggle dataset
\emph{All Lending Club Loan Data}, covering loans issued from 2007 to 2018,
\url{https://www.kaggle.com/datasets/wordsforthewise/lending-club},
accessed 2~June~2026.}
is the largest dataset considered.
Of approximately 2.26 million accepted loans, approximately 1.35 million loans
with resolved outcomes are retained. Its origination timestamps additionally
enable out-of-time evaluation.

The four datasets were selected based on public availability, documented binary
credit outcomes and variation in sample size and lending context. Public availability supports reproducible
implementation of the evaluation protocol. The datasets span three orders of
magnitude in sample size and cover different credit products, allowing the model
pipelines to be compared across heterogeneous empirical settings. Home Credit and
Lending Club additionally provide larger, industry-derived credit datasets.

\section{Model Set}\label{sec:models}

The comparison comprises logistic regression, the EBM, XGBoost and CatBoost. Logistic regression
serves as the linear transparent baseline against which any gain from added flexibility
is read. The EBM represents the non-linear intrinsically interpretable alternative,
contributing additional capacity whose shape functions remain directly inspectable.
XGBoost is the primary boosted-tree comparator, explained post hoc through TreeSHAP, and
CatBoost provides a within-family check on whether the observed patterns recur across two
boosted-tree implementations rather than defining an additional structural group. All
four models receive the shared preprocessed inputs. CatBoost's native ordered
target statistics are therefore not used. Comparisons concern the resulting
model--pipeline configurations rather than model classes in isolation.


\section{Data Preparation}\label{sec:prep}

\paragraph{Feature definition and exclusions.}
The outcome is coded so that the positive class denotes default. This operation
concerns the target label only; no supervised encoding of predictors is performed.
Identifiers, free-text fields and variables unavailable at the intended prediction
point are excluded based on the accompanying dataset documentation. Documented
sentinel and undefined values are converted to missing values according to
prespecified dataset-specific rules. For Taiwan, the payment-history variables are
retained because they describe repayment behaviour observed before the prediction
period and therefore do not constitute target leakage.

For Lending Club, platform-generated risk and pricing information, such as
\emph{grade}, \emph{sub-grade} and \emph{interest rate}, as well as
contract-derived fields such as \emph{instalment}, are excluded from the predictor
set. This prevents the models from using outputs of the platform's existing
assessment and pricing process and approximates a pre-pricing prediction point. Interest rate is stored separately only for the economic evaluation,
while origination year is used only to construct the temporal split and neither enters
the feature matrices used for model fitting or explanation.

\paragraph{Missing values and outliers.}
A prespecified missingness threshold excludes features with more than 40\%
missing observations, since their observed values cover only a limited share of
the applicant population and extensive imputation would make their representation
depend strongly on the chosen imputation rule. The resulting numbers of retained
source predictors are reported in Table~\ref{tab:datasets}. 


Retained numerical features, including ordinally encoded features, are imputed
with the median. Missing nominal values are assigned to a separate
\emph{unknown} category included in the fitted encoding schema. For nominal
variables, missingness therefore remains distinguishable from observed category
levels, and this distinction is preserved in the resulting feature-level
explanations.

Selected heavy-tailed monetary and ratio variables are capped at their
training-derived upper 99th percentiles, corresponding to upper-tail
winsorization. The affected variables are prespecified based on their documented
meaning and plausibility checks. Extreme values may represent recording
conventions, errors or genuine tail observations; the procedure does not
distinguish among these cases. The same rule is applied across all four model
pipelines before model-specific transformations.

\paragraph{Encoding and scaling.}
Ordered categorical features are mapped to consecutive integer codes following
their documented category order. For logistic regression, this representation
treats adjacent levels as equally spaced and imposes a linear effect on the
log-odds. The EBM and boosted-tree models can instead represent nonlinear changes
across the coded levels.

Nominal features are one-hot encoded without dropping a reference level. With
unseen categories represented by an all-zero indicator vector, retaining all
observed levels prevents them from being conflated with an omitted reference
category. Although the complete indicator set introduces redundancy in the
logistic-regression design matrix, estimation is stabilised by the $L_2$ penalty;
the resulting coefficients are therefore not contrasts against an omitted
baseline category.

Standardisation is fitted only within the logistic-regression pipeline. It is not
applied to the EBM or boosted-tree models, which do not require predictors to be
placed on a common scale.

\paragraph{Class imbalance.}
The non-default-to-default ratio ranges from roughly $2{:}1$ in German Credit to about
$11{:}1$ in Home Credit. Rather than resampling, imbalance is addressed by weighting the
classes inversely to their frequencies, which avoids the information loss of
under-sampling and the artefacts of synthetic over-sampling, whose effect on credit-risk
models is inconsistent \citep{zhaoResamplingTechniquesStudy2024}.
Logistic regression uses \texttt{balanced} class weights, XGBoost and CatBoost set
\texttt{scale\_pos\_weight}, and the EBM, having no class-weight parameter, is trained
with explicit balanced sample weights computed inside each fitting routine. The rule is fixed in
advance rather than tuned, so that it applies the same weighting principle through
estimator-specific interfaces.

\paragraph{Pipeline integration.}
Before pipeline fitting, preprocessing is limited to prespecified,
sample-independent operations: target coding, removal of identifiers and variables
unavailable at the prediction point, documented sentinel-value conversions, and
documentation-based ordinal mappings. Missing values, nominal category labels and
uncapped numerical values therefore remain in the stored modelling tables, which
are treated as \emph{semi-raw}.

All sample-dependent steps are
implemented within the fitted pipeline. Their parameters are estimated only from
the respective fitting sample and applied unchanged to the corresponding
validation or test data. The complete pipeline is refitted within every
cross-validation fold and once on the full training partition. Consequently, no
validation or test observation contributes to an estimated preprocessing
parameter, and the feature representation forms part of the evaluated
preprocessing-and-model pipeline.



\section{Training, Validation and Calibration}\label{sec:validation}

\paragraph{Training and test partition.}
Each dataset is partitioned once using a stratified random split, with 80\% of
the observations assigned to training and 20\% to a held-out test partition.
The same partition is used for all four model pipelines, and the observed class
proportions are approximately preserved. All model-development decisions,
including preprocessing, hyperparameter selection, probability calibration and
threshold estimation, are confined to the training partition. The test partition
is reserved for the prespecified final evaluation and explanation analyses and
does not inform any fitted parameter or selection decision.

\paragraph{Hyperparameter selection.}
Hyperparameters are selected exclusively within each training partition using
stratified five-fold cross-validation. Candidate configurations are ranked by
their mean validation ROC-AUC across the five folds.

Search spaces and budgets are prespecified separately for each model family
rather than normalised to search-space size. The logistic-regression grid is
evaluated exhaustively, whereas the EBM, XGBoost and CatBoost spaces receive
different fractional coverage under randomised search
(Table~\ref{tab:hp_all}). Exact ties in mean validation ROC-AUC are resolved in
favour of the configuration with the lower standard deviation across validation
folds.

\paragraph{Early stopping.}
For XGBoost and CatBoost, the number of boosting rounds is determined by early
stopping rather than included in the hyperparameter search space. Within each
cross-validation split, a candidate is fitted on four folds and monitored on the
fifth; training stops after 50 consecutive rounds without an improvement in
validation ROC-AUC, subject to a maximum of 2{,}000 rounds. The
median stopping point across the five folds determines the fixed round count used
for the final refit.

\paragraph{Probability calibration.}
Hyperparameter selection optimises ROC-AUC and therefore targets ranking rather
than probability calibration \citep{bequeApproachesCreditScorecard2017}. Calibration is subsequently performed entirely
within the training partition. Raw out-of-fold scores are generated through an
internal stratified three-fold cross-validation and pooled to fit an isotonic
calibrator. The selected pipeline is then refitted on the full training partition,
and the calibrator maps its raw scores to calibrated probabilities \citep{bequeApproachesCreditScorecard2017}. No test
observations are used in this procedure.

The isotonic calibrator is unweighted so that it targets the observed class
distribution. Class and sample weights are applied only to the underlying
estimator fits and are not passed to the calibration stage; doing so would instead
target a reweighted class distribution. Raw scores and calibrated probabilities
are evaluated separately according to the metric definitions in
Section~\ref{sec:eval}.

\paragraph{Reproducibility controls.}
Dataset partitions and cross-validation folds are generated from recorded seeds and
reused across all four pipelines, so the paired comparisons rest on identical splits.
A prespecified analysis plan fixes the permitted test-set analyses; test results are not
used to revise preprocessing, hyperparameters, calibration, thresholds, metric
definitions or model selection.

\section{Predictive and Economic Evaluation}\label{sec:eval}

The evaluation distinguishes ranking quality, probability accuracy,
classification at a prespecified operating point, and economic performance.
ROC-AUC and EMP serve as the primary predictive and economic measures,
respectively; the remaining metrics provide complementary diagnostics.

\paragraph{Discrimination.}
ROC-AUC is the hyperparameter-selection objective and primary discrimination
metric because it summarises how consistently the model ranks defaulters above
non-defaulters across all possible score thresholds. Average precision (AP)
complements it by summarising the precision--recall trade-off for the default
class, which is particularly relevant when defaults are less frequent than
non-defaults. Both metrics are computed from raw model scores because they
evaluate ranking rather than probability calibration. The no-skill baseline of
AP equals the default rate and therefore differs across datasets.

\paragraph{Probabilistic performance.}
Models with similar rankings can nevertheless produce substantially different
default probabilities. The Brier score is therefore computed from calibrated
probabilities to measure the mean squared error of the probability forecasts,
penalising both over- and under-confident predictions. It is the only reported
metric that directly evaluates the numerical probability estimates.

\paragraph{Threshold-dependent performance.}
Ranking metrics do not show how model scores translate into discrete
classifications at a particular operating point. Balanced accuracy and $F_1$ are
therefore reported from calibrated probabilities at a common prespecified rule.
Balanced accuracy gives equal weight to default-class recall and
non-default-class recall, preventing the majority class from dominating the
summary. $F_1$, the harmonic mean of precision and recall for the default class,
instead captures both how many defaulters are identified and how many applicants
classified as high risk actually default. It complements balanced accuracy by
focusing specifically on the positive-class decisions while ignoring true
negatives.

Applicants are classified as high risk when their predicted default probability
exceeds the default rate of the training partition. This base-rate threshold
provides a transparent, dataset-specific reference point for identifying risk
above the observed population prevalence and is fixed without test outcomes. A
threshold of $0.5$ is not used because, for calibrated probabilities, it
corresponds to an equal-error-cost rule for which no application-specific
justification is available. The base-rate threshold remains a descriptive
operating point rather than an economically optimised lending threshold; neither
balanced accuracy nor $F_1$ is interpreted as an economic decision criterion.

\paragraph{Economic performance.}
Expected Maximum Profit (Section~\ref{sec:emp}) complements the
statistical metrics by translating score-based discrimination into a normalised
economic performance measure under uncertainty about loss severity. It is used
because realised, application-specific cost and loss information is unavailable
for the public datasets. Its parameterisation is prespecified before test-set
evaluation.

The setup combines class priors estimated from the training partition with the
loss-given-default distribution specified by
\citet{verbrakenDevelopmentApplicationConsumer2014}. This distribution assigns
point masses of $0.55$ to $\lambda=0$ and $0.10$ to $\lambda=1$, with the
remaining probability mass of $0.35$ distributed uniformly over $(0,1)$. Because
the four datasets do not provide a common and consistently defined basis for
estimating product-specific loss distributions, the published specification is
used as a transparent common reference scenario. Holding it fixed across model
pipelines and datasets prevents the assumed loss distribution itself from varying
across comparisons, but does not imply that it accurately represents every
lending product. The assumed return is set to the Verbraken baseline of
$\mathrm{ROI}=0.2644$, except for Lending Club, where the mean contractual
interest rate in the training partition is used as a proxy for the return,
yielding $\mathrm{ROI}\approx0.13$.


EMP is computed directly from the raw model scores. Heterogeneous loan
exposures are not modelled. Each observation is assigned unit exposure, and no
observation-specific exposure weights are applied. EMP remains conditional on
the assumed return, loss-given-default distribution and class prevalence
\citep{xuProfitRiskdrivenCredit2024} and should not be interpreted as realised
portfolio profit. Accordingly, absolute EMP levels are interpreted only within
datasets, whereas cross-dataset comparisons concern relative model patterns. 


\section{Attribution and Agreement Measures}\label{sec:explanations}

\paragraph{Scope of the explanation analysis.}
The supplementary explanation analysis examines global source-feature rankings
produced by the four fitted model--explainer pipelines. It considers
\emph{case-sampling stability}---whether a pipeline's ranking is reproducible
when the explained cases are resampled---and \emph{cross-model agreement}---how
similar the rankings are across pipelines. Models, explainer settings and
reference data remain fixed, so the stability analysis captures sensitivity to
the composition of the explained sample rather than retraining stability.
Applicant-level explanation quality is not evaluated. The structural model
grouping is determined by model construction and does not depend on either
measure.


\paragraph{Attribution methods and reference data.}
Global rankings are constructed from additive feature contributions:
LinearSHAP for logistic regression
\citep{lundbergUnifiedApproachInterpreting2017}, the EBM's intrinsic term
contributions
\citep{louAccurateIntelligibleModels2013,
noriInterpretMLUnifiedFramework2019}, and TreeSHAP for XGBoost and CatBoost
\citep{lundbergLocalExplanationsGlobal2020}. All contributions are computed for
the final uncalibrated base estimator fitted on the complete training partition.

The procedures differ in their reference construction. LinearSHAP uses an
interventional background of 500 training observations, sampled once and held
fixed throughout the analysis. The EBM contributions follow the model's own
additive decomposition. TreeSHAP is applied in
\texttt{tree\_path\_dependent} mode and uses the path counts stored in the fitted
trees instead of a separately supplied background sample. The resulting rankings are therefore interpreted as outputs of the
complete model--explainer pipelines rather than as properties of the fitted
models alone. Because the attribution procedures use different reference and decomposition rules, their absolute magnitudes are not compared across pipelines; the analysis compares only the resulting feature rankings and top-ten sets.

\paragraph{Construction of the global rankings.}
Attributions are computed for the same held-out observations across all four
pipelines: the complete German Credit test partition and a fixed random sample
of 2{,}000 test observations from each larger dataset. They are then mapped back
from the encoded model inputs to their source features, defined as the original
predictors before one-hot encoding.

For logistic regression, XGBoost and CatBoost, the absolute contributions of all
encoded inputs belonging to a source feature are summed. The EBM adds one prior
step: each signed interaction contribution is divided equally between its two
constituent features and combined with their signed main-effect contributions,
and absolute values are taken only afterwards. Alternative allocations are possible and would alter the resulting source-feature ranking. Global importance is the mean of
these aggregated contributions across the explained cases, following the
mean-absolute attribution convention, and the features are ranked accordingly
\citep{lundbergLocalExplanationsGlobal2020}. Taking absolute values before
summing preserves contributions from category levels acting in opposing
directions.

\paragraph{Stability and agreement.}
Case-sampling stability is evaluated by recomputing the global rankings across
bootstrap samples of the explained cases, summarised by top-ten inclusion
frequency and bootstrap rank-interval width.

Cross-model agreement is evaluated within each dataset on the shared
source-feature set. Jaccard overlap quantifies the similarity of two top-ten sets
irrespective of their internal ordering; pairwise Spearman correlation
complements it by comparing the ordering of the complete source-feature
rankings \citep{nautaAnecdotalEvidenceQuantitative2023}. Both metrics are computed for
each of the six possible model pairs. Averaging these pairwise Jaccard scores
yields the mean Jaccard overlap, which serves as a single dataset-level summary
of the overall feature consensus among all four pipelines.


\section{Resampling, Stability and Uncertainty}\label{sec:comparison}

Resampling serves two distinct purposes. For predictive performance, it estimates the
sensitivity of the reported model differences to the composition of the held-out
sample. For the explanation analysis, it defines the evaluated quantity: attribution
stability measures how far a feature ranking fluctuates when the explained cases are
resampled, so the procedure constitutes the metric rather than bounding the uncertainty
around a point estimate.

\paragraph{Performance uncertainty.}
Test-set sampling uncertainty is estimated using a paired stratified bootstrap with
1{,}000 replicates per dataset \citep{raschkaModelEvaluationModel2020}. Defaults and
non-defaults are resampled separately to preserve the observed class counts. Each
bootstrap sample is used for all four pipelines, so pairwise differences are evaluated
on the same resampled applicants. All fitted components and prespecified classification
thresholds remain fixed; no model is refitted. Each metric is recomputed in every
replicate. Balanced accuracy and $F_1$ use the prespecified base-rate threshold, whereas
EMP retains its internal maximisation over candidate cut-offs.

All six metrics are reported as test-set point estimates with 95\% percentile
bootstrap intervals. Pairwise model differences are computed within each replicate, and
their intervals are obtained from the 2.5th and 97.5th percentiles of the corresponding
paired difference distributions. $\Delta$ROC-AUC and its paired interval provide the
primary summary of between-model performance differences.

\paragraph{Attribution-ranking stability.}
For each dataset, 1{,}000 paired case-bootstrap samples are drawn, with the same
resampled applicants used for all four pipelines. The fitted models, explainers and
reference data remain fixed while the global rankings are recomputed. For each of the
ten features leading a pipeline's original global ranking, its bootstrap top-ten
inclusion frequency and 95\% percentile rank interval are computed.
The reported model-level summaries are the mean inclusion frequency and mean
rank-interval width across these ten features. The two capture different things:
inclusion frequency measures set membership and reaches one when all ten features
remain in the top ten in every replicate, whereas rank-interval width, measured in
rank positions, measures movement within the ranking. High membership is therefore
compatible with substantial reordering among the ten.

\paragraph{Agreement uncertainty.}
Agreement is evaluated on the same case-bootstrap replicates. Within each, all six
pairwise Jaccard overlaps are recomputed and their mean recorded, and the reported
bounds are the corresponding percentiles of those 1{,}000 means. The resampling
therefore concerns the composition of the explained cases. The pairwise values are
never resampled among themselves. The Spearman correlations complement Jaccard and are
reported as descriptive point estimates.


\section{Sensitivity Analyses and External Validity}\label{sec:robustness}

\paragraph{Cross-dataset comparison.}
All four datasets enter the primary comparison to examine whether relative model
patterns recur across credit-scoring settings, whereas the learning-curve and
out-of-time analyses are restricted to certain datasets. Because dataset size, class
imbalance, feature representation, product context and scoring type vary jointly, cross-dataset differences are interpreted descriptively
and are not attributed to any single factor.


\paragraph{Learning curves.}
Learning curves are restricted to Home Credit and Lending Club, whose training
partitions support a broad range of sample sizes. A common size grid capped at
$250{,}000$ observations is used; grid points exceeding the available training pool
are omitted. The largest evaluated samples therefore contain $100{,}000$ observations
for Home Credit and $250{,}000$ for Lending Club. At each size, three stratified
subsamples are drawn from the fixed training partition. All four pipelines are refitted
on each subsample and evaluated on the same held-out test partition. Predictive
performance is averaged across the three refits, whereas attribution-ranking agreement
across refits is calculated as the mean pairwise top-ten Jaccard overlap over the three
refit pairs.

To make sample size the only factor that varies, model specifications are held constant
across the entire grid. One manually specified configuration per model is used: logistic
regression with $C=1$; the EBM with eight outer bags and otherwise default settings;
XGBoost with 300 boosting rounds at depth 6 and a learning rate of $0.1$; and CatBoost
with 400 iterations at the same depth and learning rate.

\paragraph{Out-of-time evaluation.}
The primary evaluation randomly partitions each pooled dataset, so its training and test
observations are drawn from the same overall sample. To examine whether the model
ordering also recurs when evaluation follows training chronologically, Lending Club, the
only dataset with a usable origination timestamp, is additionally evaluated out of time.
Candidate configurations are fitted on loans issued through 2014 and selected on the
2015 cohort. The selected pipelines are then refitted on all retained loans issued
through 2015 and evaluated on the 2016--2018 cohorts.


Origination year is used only to define the temporal partitions and is excluded from the
predictor set. The same model-specific search spaces and budgets as in the primary
comparison are retained. Throughout the Lending Club analysis, only loans with outcomes
resolved by the observation cut-off are included. Because more recent loans had less
time to reach a resolved outcome, this restriction excludes a larger share of the later
evaluation cohorts and makes the temporal estimates conditional on the subset of
matured loans. The design provides a single temporal evaluation and complements rather than replaces the random-split
benchmark used across all four datasets.

